\documentclass[10pt]{beamer}
\usetheme{Madrid}
\usecolortheme{seahorse}

% === Yellow bottom bar setup ===
\setbeamercolor{footline}{bg=yellow, fg=black}
\setbeamertemplate{footline}{%
	\leavevmode%
	\hbox{%
		\begin{beamercolorbox}[wd=\paperwidth,ht=2.5ex,dp=1ex,leftskip=1em,rightskip=1em]{footline}%
			\insertshorttitle \hfill \insertframenumber{} / \inserttotalframenumber
		\end{beamercolorbox}%
	}%
}

\setbeamertemplate{navigation symbols}{} % remove nav symbols
\setbeamertemplate{caption}[numbered]
\usepackage{graphicx}
\usepackage{amsmath,amssymb}
\usepackage{booktabs}
\usepackage{listings}
\usepackage{algorithm}
\usepackage{algpseudocode}
\usepackage{capt-of} 
\usepackage{hyperref}

\title[Beamer Lab Activity-3]{Software Tools and Technologies Laboratory\\Latex Lab Activity - 3: Beamer}
\author[Anand, Arindam]{Anand M. Baswade\and Arindam Chakraborty}
\institute[ IIT Bhilai ]{\small Faculty of Computer Science and Engineering \\ Indian Institute of Technology, Bhilai  \\ \small M.Tech Student of Computer Science and Engineering \\ Indian Institute of Technology, Bhilai}
\logo{\includegraphics[width=1.5cm]{expample-image}}

\date{15th September, 2025}
\theoremstyle{definition}
\AtBeginSection[]{
	\begin{frame}<beamer>
		\frametitle{Table of Contents}
		\tableofcontents[currentsection]
	\end{frame}
}

\begin{document}
	\begin{frame}
		\titlepage
		\vfill
	\end{frame}
	\begin{frame}{Table of Contents}
		\begin{enumerate}
			\item Introduction
			\item Why Beamer?
			\item Features
			\item Highlighting Texts
			\item Mathematical Concepts
			\item Linear Search Algorithm
			\item Images
			\item Important Equations
			\item Conclusion
			\item References
			\item Thank You
		\end{enumerate}
		\vfill
	\end{frame}
	
	\section{Introduction}
	\begin{frame}{Introduction to \LaTeX{} Beamer}
		\begin{itemize}
			\item \LaTeX{} Beamer is a powerful class for creating presentations.
			\item It allows:
			\begin{itemize}
				\item Professional-quality slides
				\item Mathematical equations
				\item Figures, tables, and references
				\item Animations and overlays
			\end{itemize}
			\item Advantages:
			\begin{itemize}
				\item Consistent formatting
				\item Easy handling of citations and references
				\item Cross-platform PDF output
			\end{itemize}
		\end{itemize}
	\end{frame}
	
	\section{Why Beamer?}
	\begin{frame}{Why Use Beamer?}
		\begin{itemize}[<+->]
			\item Beamer is designed for creating professional presentations directly in \LaTeX{}.
			\item Consistent and elegant design
			\item Full support for mathematical equations
			\item Easy management of references and citations
			\item Portable PDF output (works everywhere)
			\item Highly customizable (themes, colors, backgrounds, animations)
			\item Free and open-source — no licensing issues
		\end{itemize}
	\end{frame}
	
	\begin{frame}{Comparison with Other Tools}
		\begin{tabular}{p{0.45\textwidth} p{0.45\textwidth}}
			\toprule
			Feature & Beamer vs PowerPoint\\\midrule
			Equations & Excellent vs Limited\\
			Citations & Automatic vs Manual\\
			Portability & PDF (stable) vs File-dependent\\
			Cost & Free vs Proprietary\\
			Customization & High (code-based) vs Moderate (GUI-based)\\\bottomrule
		\end{tabular}
	\end{frame}
	
	\section{Features}
	\begin{frame}{Key Features of \LaTeX{} Beamer}
		\begin{itemize}[<+->]
			\item Predefined professional themes for elegant slides
			\item Customizable colors, fonts, and backgrounds
			\item High-quality PDF output across platforms
			\item Full support for complex mathematical equations
			\item Easy insertion of tables, figures, and graphics
			\item Step-by-step content reveal using overlays
			\item Integration with BibTeX for automatic citations
		\end{itemize}
	\end{frame}
	
	\section{Highlighting Texts}
	\begin{frame}{Blocks and Highlights Example}
		\begin{block}{Remark}
			This is a normal block. It wraps the text in a box with the same style as the presentation.
		\end{block}
		
		\begin{alertblock}{Important Theorem}
			This is an alert block. It contrasts with the regular presentation style.
		\end{alertblock}
		
		\begin{exampleblock}{Examples}
			This is an example block. Perfect for examples.
		\end{exampleblock}
	\end{frame}
	
	\section{Mathematical Concepts}
	\begin{frame}{Definition}
		\begin{definition}
			A prime number is a natural number greater than 1 that has no positive divisors other than 1 and itself.
		\end{definition}
	\end{frame}
	
	\begin{frame}{Theorem \& Lemma}
		\begin{theorem}
			\textbf{Fundamental Theorem of Arithmetic:} Every integer greater than 1 is either a prime or can be written as a unique product of prime numbers.
		\end{theorem}
		\begin{lemma}
			If a prime $p$ divides a product $ab$, then $p$ divides $a$ or $p$ divides $b$.
		\end{lemma}
	\end{frame}
	
	\section{Linear Search Algorithm}
	\begin{frame}{Linear Search Algorithm}
		\begin{algorithm}[H]
			\caption{Linear Search}
			\begin{algorithmic}[1]
				\Procedure{LinearSearch}{$A, key$}
				\For{$i=0$ to $n-1$}
				\If{$A[i]=key$}
				\State \textbf{return} $i$
				\EndIf
				\EndFor
				\State \textbf{return} -1
				\EndProcedure
			\end{algorithmic}
		\end{algorithm}
	\end{frame}
	
	\section{Corollary}
	\begin{frame}{Corollary \& Proof}
		\begin{corollary}
			The factorization of any integer into primes is unique up to the order of the factors.
		\end{corollary}
		\textbf{Proof.}
		Suppose $n$ has two different prime factorizations. By the lemma, every prime in one factorization divides one of the primes in the other factorization. This leads to a contradiction, proving uniqueness.
	\end{frame}
	
	\section{Images}
	\begin{frame}{Including an Image in a Block}
		\begin{block}{Example Image}
			\centering
			\fbox{\parbox{0.5\textwidth}{Placeholder for image}}
			\captionof{figure}{This is a sample image included in a Beamer slide.}
		\end{block}
	\end{frame}
	
	\section{Important Equations}
	\begin{frame}{Key Equations in Signal Processing}
		\begin{align*}
			X(f) &= \int_{-\infty}^{\infty} x(t) e^{-j2\pi f t}\,dt \\
			y(t) &= \int_{-\infty}^{\infty} x(\tau) h(t-\tau) \, d\tau \\
			H(f) &= \frac{Y(f)}{X(f)} = \int_{-\infty}^{\infty} h(t) e^{-j2\pi f t} \, dt
		\end{align*}
	\end{frame}
	
	\section{Conclusion}
	\begin{frame}{Conclusion}
		\begin{itemize}[<+->]
			\item \LaTeX{} Beamer is a powerful tool for creating professional, structured presentations.
			\item It supports complex mathematical equations, figures, tables, and references.
			\item Overlays and incremental reveals help guide the audience’s focus.
			\item Highly customizable themes, colors, and backgrounds make slides visually appealing.
		\end{itemize}
	\end{frame}
	
	\section{References}
	\begin{frame}[allowframebreaks]{References}
		\begin{thebibliography}{9}
			\bibitem{latexcompanion}
			Michel Goossens, Frank Mittelbach, and Alexander Samarin.
			\textit{The \LaTeX{} Companion.}
			Addison-Wesley, 1994.
			
			\bibitem{texbook}
			Donald E. Knuth.
			\textit{The TeXbook.}
			Addison-Wesley, 1984.
			
			\bibitem{kopka}
			Helmut Kopka and Patrick W. Daly.
			\textit{A Guide to LaTeX.}
			Addison-Wesley, 2003.
		\end{thebibliography}
	\end{frame}
	
	\section{Thank You}
	\begin{frame}{Thank You!}
		\centering
		\LARGE Thank You!\\[1em]
		Any Questions?
	\end{frame}
	
\end{document}